\chapter{常微分方程}\label{ch:4}
\begin{tcolorbox}[title=章首知识地图,colframe=Ink]
识别阶与线性 $\to$ 一阶方程分类 $\to$ 降阶 $\to$ 线性解空间 $\to$ 非齐次特解 $\to$ Euler、级数与积分方程。
\end{tcolorbox}
\begin{examtip}考研主线：一阶常见类型、可降阶方程和二阶常系数线性方程。常见题型：初值问题、共振特解、由积分关系建立方程。竞赛拓展：全微分、积分因子、Euler与级数解；重视存在唯一性及漏解。上述为学习任务分层，具体范围以所报考科目的当年大纲为准。\end{examtip}
\section{基本概念、分离变量与齐次方程}\label{sec:4:1}
\subsection{概念、定理与推导}

\begin{definition}{方程与解}含未知函数及其导数的关系为微分方程，出现的最高导数阶数称阶。在线性方程中未知函数及各阶导数只以一次线性组合出现，系数可依赖自变量。在一个非退化区间上，函数具有方程所需阶数的导数且逐点代入满足等式，才称为该区间上的经典解。初值问题还指定某点的函数值或若干阶导数值；满足存在唯一性条件时，这些数据确定局部唯一解。不同解区间以及仅在局部有效的隐式支路须区分。\end{definition}
\begin{theorem}{存在唯一性的常用充分条件，标准结论}若$y\prime=f(x,y)$中$f$在初值点附近连续且对$y$局部Lipschitz，则该初值问题局部存在唯一解。$f_y$连续是常用充分条件。\end{theorem}
\begin{proof}积分形式$y=y_0+\int_{x_0}^xf(t,y(t))\dd t$在充分短区间上定义压缩映射；迭代一致收敛给解，Lipschitz估计给唯一性。\end{proof}
分离方程$y\prime=p(x)q(y)$在$q(y)\ne0$区域积分$\int\frac{\dd y}{q(y)}=\int p(x)\dd x+C$；需另查$q(y)=0$的平衡解。齐次方程$y\prime=F(\frac{y}{x})$令$y=xv$，得到$xv\prime=F(v)-v$。

\begin{method}{基本概念、分离变量与齐次方程}条件：先查除去的因子是否为零。步骤：识别分离或比值齐次，积分后代初值并说明区间。易错点：一阶齐次型不等于齐次线性方程。\end{method}
\subsection{例题与方法比较}
\begin{example}{基础例题｜分离与线性}{难度：$\star$}
\label{ex:4:1}\textbf{题目} 解$y\prime=2xy,y(0)=3$。\par
\textbf{思路} 既可分离也可用线性积分因子。\par
\textbf{解答} 方法一：$\frac{\dd y}{y}=2x\dd x$，$y=C\e^{x^2}$，初值给$C=3$。方法二：$(\e^{-x^2}y)\prime=0$，故同解。零解不满足该初值，但应在通解族中保留。\par
\textbf{方法点评} 两种结构重叠，积分因子避免除以$y$。\par
\textbf{易错点或变式} 不能从分离过程直接断言零解不存在。
\end{example}
\begin{example}{考研风格例题｜齐次型}{难度：$\star\star$}
\label{ex:4:2}\textbf{题目} 解$y\prime=1+\frac{y}{x}$，$x>0$。\par
\textbf{思路} 比值代换或一阶线性。\par
\textbf{解答} 方法一：$y=xv$，$v+xv\prime=1+v$，$v=\ln x+C$，$y=x(\ln x+C)$。方法二：$y\prime-\frac{y}{x}=1$乘$\frac{1}{x}$得$(\frac{y}{x})\prime=\frac{1}{x}$。\par
\textbf{方法点评} 齐次代换解释比例结构，积分因子给最短计算。\par
\textbf{易错点或变式} 跨越$x=0$需要重新讨论方程定义。
\end{example}
\begin{example}{考研风格例题｜非唯一性}{难度：$\star\star$}
\label{ex:4:3}\textbf{题目} 初值问题$y\prime=2\sqrt{|y|},y(0)=0$是否唯一？\par
\textbf{思路} 在$y=0$处不满足局部Lipschitz。\par
\textbf{解答} 零函数为解。对任意$c\ge0$，令$y=0$（$x\le c$），$y=(x-c)^2$（$x\ge c$），连接处一阶导数为0，处处满足方程并满足初值。因此不唯一。\par
\textbf{方法点评} 缺少充分条件本身不是反证，显式构造不同解才完成判断。\par
\textbf{易错点或变式} 方程右端连续只保证局部存在，不保证唯一。
\end{example}
\section{一阶线性、伯努利与全微分}\label{sec:4:2}
\subsection{概念、定理与推导}

\begin{theorem}{一阶线性公式}设$p,q$在区间$I$上连续。对于$y\prime+p(x)y=q(x)$，令$\mu=\e^{\int p\dd x}$，则
\begin{equation}y=\mu^{-1}\left(C+\int\mu q\dd x\right).\end{equation}
\end{theorem}
\begin{proof}乘以$\mu$后左边为$(\mu y)\prime$，直接积分。伯努利方程$y\prime+py=qy^\alpha$（$\alpha\ne1$）在$y\ne0$的实值区间令$z=y^{1-\alpha}$，得到$z\prime+(1-\alpha)pz=(1-\alpha)q$。\end{proof}
\begin{theorem}{全微分与积分因子}若$M,N\in C^1$，单连通区域内$M_y=N_x$，则$M\dd x+N\dd y=\dd F$。在$N\ne0$的区域，若$\frac{M_y-N_x}{N}$为仅依赖$x$的连续函数，则$\mu(x)=\exp\int\frac{M_y-N_x}{N}\dd x$是积分因子；在$M\ne0$的区域，若$\frac{N_x-M_y}{M}$为仅依赖$y$的连续函数，同理求$\mu(y)$。\end{theorem}
局部矩形中可取$F=\int M\dd x+h(y)$，再由$F_y=N$确定$h$。积分因子公式均在分母非零的区域应用。

\begin{method}{一阶线性、伯努利与全微分}条件：先整理成规范形式。步骤：线性乘因子，伯努利降幂，全微分比交叉偏导。易错点：伯努利代换后需补查零解和实数幂定义域。\end{method}
\subsection{例题与方法比较}
\begin{example}{考研风格例题｜伯努利}{难度：$\star\star$}
\label{ex:4:4}\textbf{题目} 解$y\prime+y=y^2$。\par
\textbf{思路} 既能分离也能令$z=\frac{1}{y}$。\par
\textbf{解答} 方法一：$z\prime-z=-1$，得$z=1+C\e^x$，$y=\frac{1}{1+C\e^x}$，另有$y=0$。方法二：$\frac{\dd y}{y(y-1)}=\dd x$，部分分式积分得$\ln|\frac{y-1}{y}|=x+C_0$，化为同一非零解族，并补查$y=1,0$。\par
\textbf{方法点评} 伯努利更短，分离更容易看到两个平衡点。\par
\textbf{易错点或变式} 分母零点分隔最大存在区间。
\end{example}
\begin{example}{考研风格例题｜全微分}{难度：$\star\star$}
\label{ex:4:5}\textbf{题目} 解$(2xy+1)\dd x+(x^2+2y)\dd y=0$。\par
\textbf{思路} 交叉偏导相等。\par
\textbf{解答} 方法一：对$x$积分得$F=x^2y+x+h(y)$，由$F_y=x^2+h\prime(y)=x^2+2y$得$h=y^2$，故$x^2y+x+y^2=C$。方法二：直接识别$\dd(x^2y)+\dd x+\dd(y^2)=0$。\par
\textbf{方法点评} 势函数法系统，微分配组法适合低次多项式。\par
\textbf{易错点或变式} 积分时$h$是另一变量的函数。
\end{example}
\begin{example}{考研风格例题｜积分因子}{难度：$\star\star$}
\label{ex:4:6}\textbf{题目} 解$y\dd x+2x\dd y=0$，$x>0$。\par
\textbf{思路} 除以$2x$后线性，或给微分形式乘因子。\par
\textbf{解答} $y\prime+\frac{y}{2x}=0$，积分因子$\sqrt x$，得$y\sqrt x=C$。对原微分形式，$\frac{M_y-N_x}{N}=-\frac{1}{2x}$，积分因子$x^{-\frac{1}{2}}$；此时原式等于$2\dd(y\sqrt x)$。\par
\textbf{方法点评} 规范线性式与原微分形式的积分因子可以不同。\par
\textbf{易错点或变式} 不能忽略方程整体乘除造成的因子变化。
\end{example}
\section{降阶与常系数齐次方程}\label{sec:4:3}
\subsection{概念、定理与推导}

\begin{theorem}{线性叠加与常系数结构}线性齐次方程解空间可线性叠加。设$a,b\in\R$，$C_1,C_2\in\R$。对$y\prime\prime+ay\prime+by=0$，若特征根$r_1\ne r_2$均为实数，则$y=C_1\e^{r_1x}+C_2\e^{r_2x}$；若有实重根$r$，则$y=(C_1+C_2x)\e^{rx}$；若有非实共轭根$\alpha\pm i\beta$（$\beta>0$），则$y=\e^{\alpha x}(C_1\cos\beta x+C_2\sin\beta x)$。\end{theorem}
\begin{proof}代$y=\e^{rx}$得特征多项式。重根情形代$y=u\e^{rx}$化为$u\prime\prime=0$；复根取复指数解的实部和虚部。两基解Wronskian非零，故覆盖二阶解空间。\end{proof}
缺$y$方程令$p=y\prime(x)$；缺$x$方程在可将$p$视为$y$的函数的区间用$y\prime\prime=\frac{p\dd p}{\dd y}$。后一变换在转折点附近需另行检查。

\begin{method}{降阶与常系数齐次方程}条件：判断缺哪一个变量。步骤：缺$y$降为关于$p(x)$，缺$x$用$p(y)$；常系数写特征多项式。易错点：重根少乘$x$会使基解线性相关。\end{method}
\subsection{例题与方法比较}
\begin{example}{基础例题｜常系数}{难度：$\star$}
\label{ex:4:7}\textbf{题目} 解$y\prime\prime-3y\prime+2y=0$。\par
\textbf{思路} 特征多项式可分解。\par
\textbf{解答} $r^2-3r+2=(r-1)(r-2)$，故$y=C_1\e^x+C_2\e^{2x}$。另令$z=y\prime-y$，得$z\prime-2z=0$，再解一阶方程也得同式。\par
\textbf{方法点评} 特征根快，逐级因式分解解释算子结构。\par
\textbf{易错点或变式} 两个独立常数不可合并。
\end{example}
\begin{example}{考研风格例题｜缺因变量}{难度：$\star\star$}
\label{ex:4:8}\textbf{题目} 解$y\prime\prime=2xy\prime$。\par
\textbf{思路} 设$p=y\prime$。\par
\textbf{解答} $p\prime=2xp$给$p=C_1\e^{x^2}$，积分$y=C_2+C_1\int_0^x\e^{t^2}\dd t$。无需强求初等原函数。\par
\textbf{方法点评} 定积分表达本身就是完整解。\par
\textbf{易错点或变式} 积分下限改变只会改变常数$C_2$。
\end{example}
\begin{example}{考研风格例题｜能量降阶}{难度：$\star\star$}
\label{ex:4:9}\textbf{题目} 解$y\prime\prime=-y$并给出能量守恒。\par
\textbf{思路} 线性法给通解，乘$y\prime$给不变量。\par
\textbf{解答} 特征根$\pm i$给$y=A\cos x+B\sin x$。乘$y\prime$后$\frac\dd{\dd x}[(y\prime)^2+y^2]=0$，故能量$(y\prime)^2+y^2=A^2+B^2$守恒。\par
\textbf{方法点评} 能量关系描述相轨道，线性解给时间参数。\par
\textbf{易错点或变式} 从平方关系开根时不能全程固定一个速度符号。
\end{example}
\section{非齐次、常数变易与算子}\label{sec:4:4}

\begin{method}{待定系数试解表}
对常系数方程$P(D)y=\e^{\alpha x}P_m(x)$，取$y_p=x^s\e^{\alpha x}Q_m(x)$，$s$为$\alpha$的特征根重数。对右端$\e^{\alpha x}[P_m(x)\cos\beta x+Q_n(x)\sin\beta x]$（$\beta\ne0$），取
\[y_p=x^s\e^{\alpha x}[U_k(x)\cos\beta x+V_k(x)\sin\beta x],\quad k=\max(m,n),\]
其中$U_k,V_k$均先取完整的$k$次多项式，$s$为$\alpha+i\beta$的根重数。即使右端只有正弦，也常须同时保留余弦试项。多个右端分项求特解再相加，代回能检验每个系数。
\end{method}

\subsection{概念、定理与推导}

\begin{theorem}{非齐次结构与常数变易}设$p,q,r$在区间$I$上连续，$y_1,y_2$是$y\prime\prime+py\prime+qy=0$在$I$上的一组基本解，且$W=y_1y_2\prime-y_1\prime y_2\ne0$。非齐次方程的通解为齐次通解加一个特解，其中可取
\begin{equation}y_p=-y_1\int\dfrac{y_2r}{W}\dd x+y_2\int\dfrac{y_1r}{W}\dd x.\end{equation}
\end{theorem}
\begin{proof}设$y_p=u_1y_1+u_2y_2$，附加$u_1\prime y_1+u_2\prime y_2=0$，代回得$u_1\prime y_1\prime+u_2\prime y_2\prime=r$。解二元方程并积分。\end{proof}
待定系数法适用于多项式、指数、正余弦及其有限乘积。若激励指数为特征根且重数$s$，试解须乘$x^s$。算子移位$P(D)(\e^{\lambda x}v)=\e^{\lambda x}P(D+\lambda)v$由乘积法则得出，形式逆算子最后必须代回核验。

\begin{method}{非齐次、常数变易与算子}条件：待定系数法要求右端属于导数封闭的有限函数族。步骤：先齐次再特解，查共振，代入配系数；一般右端用常数变易。易错点：未把最高阶系数化为1就套变易公式。\end{method}
\subsection{例题与方法比较}
\begin{example}{考研风格例题｜共振}{难度：$\star\star$}
\label{ex:4:10}\textbf{题目} 解$y\prime\prime-y=\e^x$。\par
\textbf{思路} 激励指数1是简单特征根。\par
\textbf{解答} 方法一：试$y_p=Ax\e^x$，代回为$2A\e^x$，得$A=\frac{1}{2}$。方法二：设$y=\e^xv$，方程变为$v\prime\prime+2v\prime=1$，取$v=\frac{x}{2}$。通解$C_1\e^x+C_2\e^{-x}+\frac{x\e^x}{2}$。\par
\textbf{方法点评} 待定系数最快，移位法说明乘$x$的原因。\par
\textbf{易错点或变式} 试$A\e^x$只会代出零。
\end{example}
\begin{example}{考研风格例题｜三角共振}{难度：$\star\star$}
\label{ex:4:11}\textbf{题目} 解$y\prime\prime+y=\sin x$。\par
\textbf{思路} 试解需乘$x$。\par
\textbf{解答} 取$y_p=Ax\cos x$，左端$=-2A\sin x$，故$A=-\frac{1}{2}$。通解$C_1\cos x+C_2\sin x-\frac{x\cos x}{2}$。常数变易中$W=1$，代入两积分后也得到该特解加齐次项。\par
\textbf{方法点评} 特解相差齐次解不影响通解。\par
\textbf{易错点或变式} 共振特解可出现线性增长。
\end{example}
\begin{example}{考研风格例题｜一般右端}{难度：$\star\star$}
\label{ex:4:12}\textbf{题目} 解$y\prime\prime+y=\sec x$，$|x|<\frac{\pi}{2}$。\par
\textbf{思路} 右端不适合有限待定系数，使用变易。\par
\textbf{解答} 取$y_1=\cos x,y_2=\sin x,W=1$。$u_1\prime=-\tan x$给$u_1=\ln\cos x$，$u_2\prime=1$给$u_2=x$。故$y=C_1\cos x+C_2\sin x+\cos x\ln\cos x+x\sin x$。\par
\textbf{方法点评} 常数变易能处理一般连续激励。\par
\textbf{易错点或变式} 所选区间保证$\cos x>0$及右端有定义。
\end{example}
\section{Euler、幂级数与积分方程}\label{sec:4:5}
\subsection{概念、定理与推导}

\begin{theorem}{Euler方程}对$x^2y\prime\prime+axy\prime+by=f(x)$，在$x>0$令$t=\ln x,Y(t)=y(\e^t)$，则$xy\prime=Y\prime$、$x^2y\prime\prime=Y\prime\prime-Y\prime$，方程化为常系数型。\end{theorem}
\begin{proof}两次链式求导得到$Y\prime=xy\prime$与$Y\prime\prime=xy\prime+x^2y\prime\prime$。\end{proof}
\begin{extension}幂级数解在普通点把$y=\sum a_nx^n$代回，通过同次幂系数得递推。存在解析解需要系数解析且最高阶系数非零等条件。积分方程可通过求导化为微分方程，但初值与原式一致性必须回代验证。\end{extension}

\begin{method}{Euler、幂级数与积分方程}条件：Euler代换限定符号区间；级数展开须在普通点。步骤：化标准型、求系数递推或初值，再代回原式。易错点：求导会丢失常数约束。\end{method}
\subsection{例题与方法比较}
\begin{example}{竞赛风格例题｜Euler重根}{难度：$\star\star\star$}
\label{ex:4:13}\textbf{题目} 解$x^2y\prime\prime-xy\prime+y=0$，$x>0$。\par
\textbf{思路} 幂函数试解或对数换元。\par
\textbf{解答} 方法一：$y=x^m$给$(m-1)^2=0$，重根解$y=x(C_1+C_2\ln x)$。方法二：$t=\ln x$后$Y\prime\prime-2Y\prime+Y=0$，故$Y=(C_1+C_2t)\e^t$。\par
\textbf{方法点评} 幂试解简洁，对数换元解释重根中的$\ln x$。\par
\textbf{易错点或变式} 在负半轴应使用$\ln|x|$并单独建立解区间。
\end{example}
\begin{example}{竞赛风格例题｜幂级数}{难度：$\star\star\star$}
\label{ex:4:14}\textbf{题目} 用幂级数解$y\prime=y,y(0)=1$。\par
\textbf{思路} 比较相邻次数的系数。\par
\textbf{解答} 设$y=\sum_{n=0}^{\infty}a_nx^n$，有$(n+1)a_{n+1}=a_n,a_0=1$，故$a_n=\frac{1}{n!}$。所得级数对所有$x$绝对收敛且可逐项求导，解为$\e^x$。\par
\textbf{方法点评} 递推确定系数，收敛检查把形式级数变成实际解。\par
\textbf{易错点或变式} 仅求出递推并不自动证明级数收敛。
\end{example}
\begin{example}{竞赛风格例题｜积分方程}{难度：$\star\star\star$}
\label{ex:4:15}\textbf{题目} 解连续函数方程$y(x)=1+\int_0^x(x-t)y(t)\dd t$。\par
\textbf{思路} 连续核使右端可微两次。\par
\textbf{解答} 求导得$y\prime=\int_0^xy(t)\dd t$，再求导$y\prime\prime=y$。原式给$y(0)=1,y\prime(0)=0$，故$y=\cosh x$。回代积分为$\cosh x-1$，原式成立。\par
\textbf{方法点评} 微分方程与两条初始约束共同恢复原积分关系。\par
\textbf{易错点或变式} 只写$y\prime\prime=y$会多出不合法的解。
\end{example}
\begin{warning}求解过程中除以$y$、$y\prime$或$x$都可能丢解或缩小区间；特解不要求唯一；方程通解必须有相应数量的独立常数并经过代回检验。\end{warning}


\section{一阶隐式方程与微分方程基本定理}\label{mat:ode_theory}
\begin{examtip}补入线素场、隐式类型、奇解、延展、连续依赖与参数变分。\end{examtip}
\subsection{线素场与Euler折线}\label{mat:ode_theory:direction}
\begin{method}{线素场与Euler折线}\textbf{适用条件} $y\prime=f(x,y)$连续，在计算区间适当光滑。\par\textbf{实施与依据} 在网格点用斜率$f(x_n,y_n)$作线段，迭代$y_{n+1}=y_n+h f(x_n,y_n)$。\end{method}
\begin{example}{专题例题｜线素场与Euler折线}{难度：$\star\star$}
\textbf{题目} 对$y\prime=y,y(0)=1$使用Euler法到$t=1$。\par\textbf{思路与解答} 取$h=\frac{1}{N}$，得到$y_N=(1+\frac{1}{N})^N\to\e$。若$f$对$y$为Lipschitz且局部截断误差为$O(h^2)$，离散Gronwall给全局误差$O(h)$。\par\textbf{方法点评与边界} 存在唯一性与数值稳定性是不同问题，步长过大可产生伪振荡。
\end{example}
\subsection{Riccati方程的已知特解降阶}\label{mat:ode_theory:riccati}
\begin{method}{Riccati方程的已知特解降阶}\textbf{适用条件} $y\prime=a(x)y^2+b(x)y+c(x)$且已知特解$y_p$。\par\textbf{实施与依据} 在$y\ne y_p$区间令$y=y_p+\frac{1}{u}$，化为$u\prime+(2ay_p+b)u=-a$。\end{method}
\begin{example}{专题例题｜Riccati方程的已知特解降阶}{难度：$\star\star$}
\textbf{题目} 解$y\prime=y^2-1$并说明降阶。\par\textbf{思路与解答} 已知$y_p=1$，得$u\prime+2u=-1$，$u=C\e^{-2x}-\frac{1}{2}$，所以$y=1+\frac{1}{C\e^{-2x}-\frac{1}{2}}$，另补$y=1$。也可分离积分$\frac12\ln|\frac{y-1}{y+1}|=x+C_0$并补$\pm1$。\par\textbf{方法点评与边界} 代换可能遗漏已知特解，必须补回。
\end{example}
\subsection{Clairaut与d’Alembert型}\label{mat:ode_theory:clairaut}
\begin{method}{Clairaut与d’Alembert型}\textbf{适用条件} 含$p=y\prime$的一阶隐式方程。\par\textbf{实施与依据} Clairaut型$y=xp+\phi(p)$求导给$(x+\phi\prime(p))p\prime=0$，分别求常$p$直线族与包络。\end{method}
\begin{example}{专题例题｜Clairaut与d’Alembert型}{难度：$\star\star$}
\textbf{题目} 解$y=xy\prime+(y\prime)^2$。\par\textbf{思路与解答} 常$p=C$给直线族$y=Cx+C^2$；另一支$x+2p=0$给$p=-\frac{x}{2}$，代回$y=-\frac{x^2}{4}$，并验证其导数为$-\frac{x}{2}$。这是奇解。\par\textbf{方法点评与边界} 只保留$p\prime=0$会漏掉包络奇解；对$p$求导的推导先在二阶可导的支路进行。若允许仅$C^1$的解，还可在直线与包络的相切点拼接，须代回原式检查，不能把两个光滑解族直接宣称为全部$C^1$解。
\end{example}
\subsection{正交轨线与模型}\label{mat:ode_theory:trajectory}
\begin{method}{正交轨线与模型}\textbf{适用条件} 给定曲线族且斜率有限非零的局部区域。\par\textbf{实施与依据} 消去参数得到原斜率$m(x,y)$，正交族满足新斜率$-\frac{1}{m}$。\end{method}
\begin{example}{专题例题｜正交轨线与模型}{难度：$\star\star$}
\textbf{题目} 求抛物线族$y=Cx^2$的正交轨线。\par\textbf{思路与解答} 原斜率$y\prime=\frac{2y}{x}$，正交族$y\prime=-\frac{x}{2y}$，积分$2y\dd y=-x\dd x$得$y^2+\frac{x^2}{2}=C_1$。退化点需用切向几何补查。\par\textbf{方法点评与边界} 正交法在斜率0或无穷处不直接相除。
\end{example}
\subsection{Picard迭代与存在唯一性证明}\label{mat:ode_theory:picard}
\begin{method}{Picard迭代与存在唯一性证明}\textbf{适用条件} $f$在矩形内连续并对$y$统一Lipschitz。\par\textbf{实施与依据} 令$y_0$为常函数，$y_{n+1}=y_0+\int_{t_0}^tf(s,y_n(s))\dd s$，用阶乘误差或压缩估计证明一致收敛。\end{method}
\begin{example}{专题例题｜Picard迭代与存在唯一性证明}{难度：$\star\star$}
\textbf{题目} 从$y\prime=y,y(0)=1$的迭代得到解。\par\textbf{思路与解答} $y_0=1,y_1=1+t,y_2=1+t+\frac{t^2}{2}$，一般$y_n=\sum_{k=0}^n\frac{t^k}{k!}$。在有限区间尾项一致趋零，极限为$\e^t$并满足积分方程。\par\textbf{方法点评与边界} 迭代的收敛与极限满足方程必须各自证明。
\end{example}
\subsection{最大解区间与爆破}\label{mat:ode_theory:extension}
\begin{method}{最大解区间与爆破}\textbf{适用条件} 局部存在唯一解。\par\textbf{实施与依据} 若解图像在有限端点前始终留在定义域的一个紧子集，局部定理允许继续延展；不能延展时可能逼近边界或无界。\end{method}
\begin{example}{专题例题｜最大解区间与爆破}{难度：$\star\star$}
\textbf{题目} 求$y\prime=y^2,y(0)=1$的最大区间。\par\textbf{思路与解答} 解$y=\frac{1}{1-t}$，在$t=1$爆破，含0的最大存在区间为$(-\infty,1)$。不能跨过1把同一公式称为同一初值解的延拓。\par\textbf{方法点评与边界} 公式在其他区间有定义不代表能穿过奇点连续延拓。
\end{example}
\subsection{连续依赖与Gronwall}\label{mat:ode_theory:dependence}
\begin{method}{连续依赖与Gronwall}\textbf{适用条件} 同一ODE，右端对$y$有Lipschitz常数$L$。\par\textbf{实施与依据} 两个解的差满足$|u-v|\le|u_0-v_0|+L\int|u-v|$，使用Gronwall。\end{method}
\begin{example}{专题例题｜连续依赖与Gronwall}{难度：$\star\star$}
\textbf{题目} 给出初值扰动的误差界。\par\textbf{思路与解答} 在共同区间$t\ge t_0$有$|u(t)-v(t)|\le|u_0-v_0|\e^{L(t-t_0)}$。右端再有一致扰动$\delta$时，误差上界增加$\frac{\delta(\e^{L(t-t_0)}-1)}{L}$，$L=0$按极限解释。\par\textbf{方法点评与边界} 该界说明有限时间稳定，不自动给无穷时间稳定。
\end{example}
\subsection{对初值与参数的可微依赖}\label{mat:ode_theory:variation}
\begin{method}{对初值与参数的可微依赖}\textbf{适用条件} $f$对状态与参数连续可微。\par\textbf{实施与依据} 对积分方程求参数导数，得到线性变分方程。\end{method}
\begin{example}{专题例题｜对初值与参数的可微依赖}{难度：$\star\star$}
\textbf{题目} $y\prime=ay,y(0)=y_0$，求对$a$和$y_0$的敏感度。\par\textbf{思路与解答} $y=y_0\e^{at}$，故$\partial_a y=t y_0\e^{at}$、$\partial_{y_0}y=\e^{at}$。一般敏感度$z$满足$z\prime=f_yz+f_a$并配相应初值。\par\textbf{方法点评与边界} 需要正则性才能对解族求参数导数。
\end{example}

\section{高阶方程、线性系统与稳定性}\label{mat:ode_systems}
\begin{examtip}本节讨论高阶理论、Laplace变换、振动、矩阵解法、相平面与Lyapunov稳定性。\end{examtip}
\subsection{Wronskian与Abel公式}\label{mat:ode_systems:wronskian}
\begin{method}{Wronskian与Abel公式}\textbf{适用条件} 线性$n$阶方程最高阶系数非零，系数连续。\par\textbf{实施与依据} 基本解的Wronskian非零；标准型中$W\prime=-a_{n-1}W$。\end{method}
\begin{example}{专题例题｜Wronskian与Abel公式}{难度：$\star\star$}
\textbf{题目} 验证$y_1=\e^x,y_2=x\e^x$为$(D-1)^2y=0$的基本解。\par\textbf{思路与解答} $W=\e^{2x}\ne0$。方程$y\prime\prime-2y\prime+y=0$的Abel公式给$W\prime=2W$，与实际值一致。\par\textbf{方法点评与边界} 对任意函数组，Wronskian为零不总等价于线性相关；线性ODE解族的条件不能省。
\end{example}
\subsection{Laplace解初值问题}\label{mat:ode_systems:laplaceode}
\begin{method}{Laplace解初值问题}\textbf{适用条件} 半轴函数指数阶，变换存在。\par\textbf{实施与依据} $\mathcal L(y\prime)=sY-y(0)$，二阶再减初始导数；把ODE化为代数式再逆变换。\end{method}
\begin{example}{专题例题｜Laplace解初值问题}{难度：$\star\star$}
\textbf{题目} 解$y\prime\prime+y=1,y(0)=y\prime(0)=0$。\par\textbf{思路与解答} 变换得$(s^2+1)Y=\frac{1}{s}$，$Y=\frac{1}{s}-\frac{s}{s^2+1}$，故$y=1-\cos t$。代回初值与方程均成立。\par\textbf{方法点评与边界} 不连续激励可用阶跃函数移位，冲激需分布框架另行解释。
\end{example}
\subsection{阻尼、共振与强迫振动}\label{mat:ode_systems:oscillation}
\begin{method}{阻尼、共振与强迫振动}\textbf{适用条件} $my\prime\prime+cy\prime+ky=F_0\cos\omega t$，$m,k>0,c\ge0,\omega\ge0$。\par\textbf{实施与依据} 由特征根分欠阻尼、临界与过阻尼。当$(k-m\omega^2)^2+c^2\omega^2>0$时，周期特解的幅值为$\frac{|F_0|}{\sqrt{(k-m\omega^2)^2+c^2\omega^2}}$；若$c=0,\omega=\sqrt{k/m},F_0\ne0$，发生无阻尼共振，不存在有限幅值的周期特解。\end{method}
\begin{example}{专题例题｜阻尼、共振与强迫振动}{难度：$\star\star$}
\textbf{题目} 解释无阻尼共振。\par\textbf{思路与解答} $c=0,\omega=\sqrt{\frac{k}{m}}$时激励频率是齐次特征频率，特解含$t\sin\omega t$，振幅线性增长；待定系数不能只取常幅三角式。\par\textbf{方法点评与边界} 有阻尼时幅值公式分母不为零，峰值频率不必恰等于固有频率。
\end{example}
\subsection{基本解矩阵与常数变易}\label{mat:ode_systems:fundamental}
\begin{method}{基本解矩阵与常数变易}\textbf{适用条件} $\bm x\prime=A(t)\bm x+\bm f(t)$。\par\textbf{实施与依据} 若$\Phi\prime=A\Phi,\det\Phi\ne0$，则$\bm x=\Phi[\bm c+\int\Phi^{-1}\bm f]$。\end{method}
\begin{example}{专题例题｜基本解矩阵与常数变易}{难度：$\star\star$}
\textbf{题目} 解$\bm x\prime=\operatorname{diag}(1,-1)\bm x$。\par\textbf{思路与解答} $\Phi=\operatorname{diag}(\e^t,\e^{-t})$，故$x=C_1\e^t,y=C_2\e^{-t}$。行列式1与$\exp\int\operatorname{tr}A=1$一致。\par\textbf{方法点评与边界} 矩阵乘法顺序不能交换，变系数矩阵一般也不在不同时刻对易。
\end{example}
\subsection{矩阵指数与Jordan块}\label{mat:ode_systems:jordan}
\begin{method}{矩阵指数与Jordan块}\textbf{适用条件} 常系数线性系统。\par\textbf{实施与依据} $\e^{At}=\sum \frac{A^nt^n}{n!}$；若$A=\lambda I+N$、$N$幂零，指数拆成有限多项式乘指数。\end{method}
\begin{example}{专题例题｜矩阵指数与Jordan块}{难度：$\star\star$}
\textbf{题目} 求$A=\begin{pmatrix}1&1\\0&1\end{pmatrix}$的矩阵指数。\par\textbf{思路与解答} $N^2=0$，故$\e^{At}=\e^t(I+tN)=\e^t\begin{pmatrix}1&t\\0&1\end{pmatrix}$。\par\textbf{方法点评与边界} 重特征值不能一律按可对角化处理。
\end{example}
\subsection{相平面、节点、鞍点与焦点}\label{mat:ode_systems:phase}
\begin{method}{相平面、节点、鞍点与焦点}\textbf{适用条件} 二维自治系统在平衡点附近。\par\textbf{实施与依据} 线性系统按特征值实部与实虚性分类；非线性在双曲平衡点可用线性化判断局部稳定类型。\end{method}
\begin{example}{专题例题｜相平面、节点、鞍点与焦点}{难度：$\star\star$}
\textbf{题目} 分析$x\prime=x,y\prime=-y$。\par\textbf{思路与解答} 轨道满足$xy=C$，一方向指数增、一方向指数减，原点为不稳定鞍点；$x=0$为稳定方向，$y=0$为不稳定方向。\par\textbf{方法点评与边界} 零实部特征值属于临界情形，线性化不能直接定稳定性。
\end{example}
\subsection{Lyapunov稳定与渐近稳定}\label{mat:ode_systems:lyapunov}
\begin{method}{Lyapunov稳定与渐近稳定}\textbf{适用条件} 平衡点附近有正定可微函数$V$。\par\textbf{实施与依据} $\dot V=\nabla V\cdot\bm f\le0$给出稳定性的常用充分条件；严格负定并有适当局部条件给渐近稳定。\end{method}
\begin{example}{专题例题｜Lyapunov稳定与渐近稳定}{难度：$\star\star$}
\textbf{题目} 分析$x\prime=-x-y,y\prime=x-y$的原点。\par\textbf{思路与解答} 取$V=\frac{x^2+y^2}{2}$，则$\dot V=x(-x-y)+y(x-y)=-x^2-y^2=-2V$，所以$V(t)=V(0)\e^{-2t}$，所有解趋向原点，且初始小扰动始终保持小。因此原点全局渐近稳定。\par\textbf{方法点评与边界} 候选函数的符号、沿解导数及适用邻域必须逐项验证。
\end{example}
\subsection{极限环与径向方程}\label{mat:ode_systems:cycle}
\begin{method}{极限环与径向方程}\textbf{适用条件} 平面自治系统可转极坐标。\par\textbf{实施与依据} 由$r\prime=\frac{xx\prime+yy\prime}{r}$和$\theta\prime=\frac{xy\prime-yx\prime}{r^2}$分析。\end{method}
\begin{example}{专题例题｜极限环与径向方程}{难度：$\star\star$}
\textbf{题目} $x\prime=x(1-r^2)-y,y\prime=y(1-r^2)+x$，分析轨道。\par\textbf{思路与解答} $r\prime=r(1-r^2),\theta\prime=1$。$0<r<1$时半径增，$r>1$时减，故$r=1$为稳定周期轨道即稳定极限环；原点不稳定。\par\textbf{方法点评与边界} 周期轨道不一定孤立；中心附近一族闭轨道不是极限环。
\end{example}

\section{差分方程、首次积分与一阶偏微分方程}\label{mat:difference_pde}
\begin{examtip}本节通过递推、守恒量和特征线法，建立离散系统与一阶PDE的计算入口。\end{examtip}
\subsection{一阶线性差分}\label{mat:difference_pde:difference1}
\begin{method}{一阶线性差分}\textbf{适用条件} $x_{n+1}=ax_n+b_n$。\par\textbf{实施与依据} 逐步迭代得$x_n=a^nx_0+\sum_{k=0}^{n-1}a^{n-1-k}b_k$；常数激励找不动点。\end{method}
\begin{example}{专题例题｜一阶线性差分}{难度：$\star\star$}
\textbf{题目} 解$x_{n+1}=\frac{x_n}{2}+1,x_0=0$。\par\textbf{思路与解答} 不动点为2，令$z_n=x_n-2$，则$z_{n+1}=\frac{z_n}{2}$，所以$x_n=2-2^{1-n}\to2$。\par\textbf{方法点评与边界} 当$|a|\ge1$时不动点存在不代表迭代收敛。
\end{example}
\subsection{二阶常系数差分}\label{mat:difference_pde:difference2}
\begin{method}{二阶常系数差分}\textbf{适用条件} $x_{n+2}+px_{n+1}+qx_n=f_n$。\par\textbf{实施与依据} 齐次用特征方程，重根乘$n$；非齐次用待定系数与共振修正。\end{method}
\begin{example}{专题例题｜二阶常系数差分}{难度：$\star\star$}
\textbf{题目} 解Fibonacci递推$F_{n+2}=F_{n+1}+F_n,F_0=0,F_1=1$。\par\textbf{思路与解答} 特征根$\phi=\frac{1+\sqrt5}{2},\psi=\frac{1-\sqrt5}{2}$，解$F_n=\frac{\phi^n-\psi^n}{\sqrt5}$，因$|\psi|<1$得$\frac{F_{n+1}}{F_n}\to\phi$。\par\textbf{方法点评与边界} 离散指数$r^n$与连续指数$\e^{rx}$形式相似，但特征多项式的形成机制不同。
\end{example}
\subsection{非线性迭代与局部稳定}\label{mat:difference_pde:nonlinear}
\begin{method}{非线性迭代与局部稳定}\textbf{适用条件} $x_{n+1}=T(x_n)$，$T(L)=L$。\par\textbf{实施与依据} $|T\prime(L)|<1$给邻域吸引；大于1给局部排斥，等于1时需高阶分析。\end{method}
\begin{example}{专题例题｜非线性迭代与局部稳定}{难度：$\star\star$}
\textbf{题目} 分析$T(x)=rx(1-x)$的非零不动点。\par\textbf{思路与解答} $L=1-\frac{1}{r}$，$T\prime(L)=2-r$。当$1<r<3$时模小于1，非零不动点局部吸引；临界$r=3$不能只凭此判定后续动力行为。\par\textbf{方法点评与边界} 局部稳定不表示所有初值都收敛到该点。
\end{example}
\subsection{首次积分与降维}\label{mat:difference_pde:firstintegral}
\begin{method}{首次积分与降维}\textbf{适用条件} 自治向量场$\bm x\prime=\bm f(\bm x)$。\par\textbf{实施与依据} $I$是首次积分当且仅当沿轨道$\nabla I\cdot\bm f=0$。\end{method}
\begin{example}{专题例题｜首次积分与降维}{难度：$\star\star$}
\textbf{题目} 系统$x\prime=y,y\prime=-x$的首次积分。\par\textbf{思路与解答} $I=x^2+y^2$，$\dot I=2xy-2xy=0$，轨道位于圆上；结合$\theta\prime=-1$恢复运动方向。\par\textbf{方法点评与边界} 一个首次积分只给轨道约束，不一定给完整时间参数解。
\end{example}
\subsection{线性齐次PDE与特征线}\label{mat:difference_pde:transport}
\begin{method}{线性齐次PDE与特征线}\textbf{适用条件} $a(x,y)u_x+b(x,y)u_y=0$且特征方向非零。\par\textbf{实施与依据} 解$\frac{\dd x}{\dd s}=a,\frac{\dd y}{\dd s}=b,\frac{\dd u}{\dd s}=0$，以特征不变量表达任意函数。\end{method}
\begin{example}{专题例题｜线性齐次PDE与特征线}{难度：$\star\star$}
\textbf{题目} 解$u_x+2u_y=0$。\par\textbf{思路与解答} 特征线$y-2x=C$，沿线$u$不变，通解$u=\Phi(y-2x)$；直接求偏导代入为0。\par\textbf{方法点评与边界} 初始曲线须非特征才能一般地给局部唯一解。
\end{example}
\subsection{拟线性PDE与特征相交}\label{mat:difference_pde:quasilinear}
\begin{method}{拟线性PDE与特征相交}\textbf{适用条件} $a(x,y,u)u_x+b(x,y,u)u_y=c(x,y,u)$。\par\textbf{实施与依据} 同步求解三维特征系统$x\prime=a,y\prime=b,u\prime=c$。\end{method}
\begin{example}{专题例题｜拟线性PDE与特征相交}{难度：$\star\star$}
\textbf{题目} 解$u_t+u u_x=0,u(x,0)=g(x)$的局部隐式形式。\par\textbf{思路与解答} 从初始点$\xi$出发，$u=g(\xi)$保持不变，$x=\xi+t g(\xi)$。若$1+t g\prime(\xi)\ne0$可反解$\xi$，从而$u(x,t)=g(\xi(x,t))$。特征相交时经典单值光滑解可能失效。\par\textbf{方法点评与边界} 不能把越过特征交叉后的隐式多值表达当作经典解。
\end{example}




\section{常微分方程逐类深化与多方法比较}\label{deep:ode_source_deep}

\begin{examtip}按方程结构、解的理论、线性算子和方程组逐层展开，区分通解、特解、奇解与最大存在区间。\end{examtip}

\subsection{可分离、齐次与降阶中的失解}\label{deep:ode_source_deep:lost}

\begin{method}{可分离、齐次与降阶中的失解}\textbf{适用条件} 除以含未知函数的因子前，先检查该因子为零的解。\par\textbf{步骤与依据} 列出平衡解；在非零支路分离；最后用初值和唯一性确定可拼接性。\par\textbf{易错点} 隐式式子满足不等于所有分支都满足原微分方程。\end{method}

\begin{example}{考研风格｜分离变量与线性化双解}{难度：$\star\star\star$}\label{deep:ex:ode_source_deep:logistic}\textbf{题目} 解$y'=y(1-y),y(0)=y_0$，并说明平衡解。\par\textbf{思路} 分离变量，或对非零解令$z=\frac{1}{y}$。\par\textbf{详细解答} 平衡解为$y=0,1$。非平衡解$\int[\frac1y+\frac1{1-y}]\dd y=\int\dd x$，得$\ln|\frac{y}{1-y}|=x+C$，于是$y(x)=\frac{y_0\e^x}{1+y_0(\e^x-1)}$，在分母不为0且含0的最大区间成立。第二法$z'=-\frac{y'}{y^2}=1-z$，所以$z=1+Ce^{-x}$，回代同式。\par\textbf{方法点评} 两法均需补$y=0$；$y=1$已包含在公式中。\par\textbf{变式与易错点} 若$y_0<0$或$y_0>1$，分母可能在有限时间为0，不能宣称全实轴解。\end{example}

\begin{example}{考研风格｜齐次方程与几何不变量}{难度：$\star\star$}\label{deep:ex:ode_source_deep:homogeneous}\textbf{题目} 解$y'=\frac{x+y}{x}$，$x\ne0$。\par\textbf{思路} 这是齐次形式，也是一阶线性式。\par\textbf{详细解答} \textbf{齐次换元：}$y=xv$，$y'=v+xv'$，所以$xv'=1$，$v=\ln|x|+C$，故$y=x(\ln|x|+C)$。\par\textbf{线性法：}$y'-\frac{y}{x}=1$，积分因子在固定半轴可取$\frac{1}{x}$，得到$(\frac{y}{x})'=\frac{1}{x}$，同解。\par\textbf{方法点评} 同一方程有多种分类时选微分最简单者。\par\textbf{变式与易错点} 解在正、负半轴分别取常数；不能经过系数奇点0。\end{example}

\begin{example}{考研风格｜Bernoulli转线性}{难度：$\star\star\star$}\label{deep:ex:ode_source_deep:bernoulli}\textbf{题目} 解$y'+y=xy^2$。\par\textbf{思路} 对非零解令$z=y^{-1}$。\par\textbf{详细解答} $z'=-\frac{y'}{y^2}$，原式除以$y^2$得$-z'+z=x$，即$z'-z=-x$。取积分因子$\e^{-x}$，有$(z\e^{-x})'=-x\e^{-x}$，积分得$z=x+1+C\e^x$，所以$y=\frac{1}{x+1+C\e^x}$，另有$y=0$。可将非零解视为Riccati在已知特解0下的线性化。\par\textbf{方法点评} 变换由幂次$1-n$决定，负号最易出错。\par\textbf{变式与易错点} 通解只在分母不为0的连通区间成立。\end{example}

\begin{example}{竞赛风格｜Riccati已知特解法}{难度：$\star\star\star$}\label{deep:ex:ode_source_deep:riccati}\textbf{题目} 解$y'=y^2-2xy+x^2+1$。\par\textbf{思路} 先发现特解$y_p=x$，再令$y=x+z$。\par\textbf{详细解答} $1+z'=z^2+1$，所以$z'=z^2$。$z=0$给$y=x$；其余$-\frac{1}{z}=x+C$，得$y=x-\frac{1}{x+C}$。第二法令$z=-\frac{u'}{u}$，则$z'=z^2$转为$u''=0$，$u=Ax+B$，回代得同一族。\par\textbf{方法点评} 已知Riccati特解可以将一般非线性式降成Bernoulli或二阶线性式。\par\textbf{变式与易错点} 线性化中的$u$不能为0；$u$乘非零常数不改变$z$，自由常数数目仍为一个。\end{example}

\subsection{全微分、积分因子与轨线}\label{deep:ode_source_deep:exact}

\begin{method}{全微分、积分因子与轨线}\textbf{适用条件} $M\dd x+N\dd y=0$，在单连通区域内系数有连续偏导。\par\textbf{步骤与依据} 先检$M_y=N_x$；若不成立，试$\frac{M_y-N_x}{N}$仅依赖$x$，或$\frac{N_x-M_y}{M}$仅依赖$y$。\par\textbf{易错点} 局部恰当与全局势函数有拓扑差别；被除分母为0的点须独立检查。\end{method}

\begin{example}{考研风格｜单变量积分因子}{难度：$\star\star\star$}\label{deep:ex:ode_source_deep:factor}\textbf{题目} 解$(2xy+y)\dd x+x\dd y=0$，$x\ne0$。\par\textbf{思路} 检测恰当性，找只依赖$x$的积分因子。\par\textbf{详细解答} $M_y=2x+1,N_x=1$，故$\frac{M_y-N_x}{N}=2$，取$\mu(x)=\e^{2x}$。势函数$F=x y\e^{2x}$，因$F_x=(2x+1)y\e^{2x},F_y=x\e^{2x}$，解为$xy\e^{2x}=C$。第二法直接写$y'+(2+\frac{1}{x})y=0$，分离得$y=\frac{C\e^{-2x}}{x}$。\par\textbf{方法点评} 积分因子与一阶线性法本质相同，但势函数形式更易检查。\par\textbf{变式与易错点} $x=0$不是允许的自变量区间，不能由隐式式子擅自扩充。\end{example}

\begin{example}{竞赛风格｜正交轨线与等角轨线}{难度：$\star\star\star$}\label{deep:ex:ode_source_deep:orthogonal}\textbf{题目} 求圆族$x^2+y^2=C$的正交轨线。\par\textbf{思路} 先消参数求斜率，再取负倒数。\par\textbf{详细解答} 圆族斜率$y'=-\frac{x}{y}$，正交轨线在非垂直点满足$y'=\frac{y}{x}$，故$y=kx$。竖直直线$x=0$也为轨线，由几何补入。若与原族夹角为$\alpha$，新斜率$m$满足$\frac{m-m_0}{1+mm_0}=\tan\alpha$，解得$m=\frac{m_0+\tan\alpha}{1-m_0\tan\alpha}$，用于一般等角轨线。\par\textbf{方法点评} 正交轨线的参数是曲线族参数，不能与原圆半径混淆。\par\textbf{变式与易错点} 负倒数规则在水平、竖直点需按切向量理解。\end{example}

\subsection{一阶隐式方程、参数法与奇解}\label{deep:ode_source_deep:implicit}

\begin{method}{一阶隐式方程、参数法与奇解}\textbf{适用条件} 方程可写为$F(x,y,p)=0,p=y'$，且局部可对一个变量求解。\par\textbf{步骤与依据} Clairaut型先微分拆分支；Lagrange型以$p$为参数，解关于$x(p)$的线性方程。\par\textbf{易错点} 微分后得到的解须代回原式，可能出现奇解或丢失常斜率解。\end{method}

\begin{example}{竞赛风格｜Clairaut包络的两条推导}{难度：$\star\star\star$}\label{deep:ex:ode_source_deep:clairaut}\textbf{题目} 解$y=xy'+(y')^2$并求奇解。\par\textbf{思路} 令$p=y'$，微分后分成$p'=0$或$x+2p=0$。\par\textbf{详细解答} 微分给$p=p+(x+2p)p'$。$p'=0$时$p=C$，通解$y=Cx+C^2$。另一支$p=-\frac{x}{2}$，代回得$y=-\frac{x^2}{4}$，导数确为$-\frac{x}{2}$，是奇解。第二法把直线族视为包络，由$\partial_C(Cx+C^2)=x+2C=0$消去$C$，同得抛物线。\par\textbf{方法点评} 微分法验证方程，包络法说明几何构造。\par\textbf{变式与易错点} 奇解不等同于给通解选一个常数。上述微分拆支先用于二阶可导的局部支路；仅要求$C^1$时，直线与抛物线可在相切点拼接，仍须回代原方程验证。\end{example}

\begin{example}{竞赛风格｜Lagrange参数解}{难度：$\star\star\star$}\label{deep:ex:ode_source_deep:lagrange}\textbf{题目} 解$y=2xy'+(y')^2$。\par\textbf{思路} 令$p=y'$，将$x$作为$p$的函数。\par\textbf{详细解答} 微分得$p=2p+(2x+2p)p'$，即$-\frac{p\dd x}{\dd p}=2x+2p$。当$p\ne0$，有$\frac{\dd x}{\dd p}+\frac{2x}{p}=-2$，乘$p^2$得$(p^2x)'=-2p^2$，所以$x=Cp^{-2}-\frac{2p}{3}$，$y=2xp+p^2=\frac{2C}{p}-\frac{p^2}{3}$。另$p=0$代原式得$y=0$，是需补的常斜率解。参数解在$\frac{\dd x}{\dd p}\ne0$段可表成$y(x)$。\par\textbf{方法点评} 参数解本身就是完整答案，不必强行消参数。\par\textbf{变式与易错点} 把$p=0$除掉前必须单列检查。\end{example}

\subsection{存在唯一、延拓与初值依赖}\label{deep:ode_source_deep:exist}

\begin{method}{存在唯一、延拓与初值依赖}\textbf{适用条件} $f$连续保证局部存在，局部Lipschitz保证唯一。\par\textbf{步骤与依据} 积分方程写出Picard迭代；用压缩或阶乘误差证明收敛；延拓时检查解是否离开紧集。\par\textbf{易错点} 存在、唯一、全局存在是三个不同结论。\end{method}

\begin{theorem}{Picard迭代的局部构造}设$a,b>0$，$f$在矩形$|t-t_0|\le a,|y-y_0|\le b$内连续，且$|f|\le M$、$|f(t,u)-f(t,v)|\le L|u-v|$，$M,L\ge0$。取$0<h\le\min(a,\frac{b}{M})$且$Lh<1$（$M=0$时将$\frac{b}{M}$解释为$+\infty$），则初值问题在$[t_0-h,t_0+h]$有唯一解。\end{theorem}\begin{proof}在连续函数的闭球$\|y-y_0\|_\infty\le b$上定义$Ty=y_0+\int_{t_0}^tf(s,y(s))\dd s$。有$\|Ty-y_0\|_\infty\le Mh\le b$，且$\|Tu-Tv\|_\infty\le Lh\|u-v\|_\infty$，故压缩映射给唯一不动点。积分方程求导即原方程。较长局部区间可分段延拓。\end{proof}

\begin{example}{竞赛风格｜连续而不唯一}{难度：$\star\star\star$}\label{deep:ex:ode_source_deep:nonunique}\textbf{题目} 讨论$y'=2\sqrt{|y|},y(0)=0$的非负解。\par\textbf{思路} 原点不满足对$y$局部Lipschitz，允许等待后离开零解。\par\textbf{详细解答} 对任意$c\ge0$，$y_c(t)=0$（$t\le c$）、$y_c(t)=(t-c)^2$（$t\ge c$）是$C^1$解，连接点左右导数均0，满足方程及初值。还有恒零解。因此连续性不能保证唯一。若限定非负解，正值段由分离$\frac{\dd y}{\sqrt y}=2\dd t$给平方支路；零值段可与其如上连接。\par\textbf{方法点评} 直接验连接点比口头说“可拼接”更严谨。\par\textbf{变式与易错点} 方程$y'=y$满足Lipschitz，不能作同样等待拼接。\end{example}

\begin{example}{竞赛风格｜最大存在区间}{难度：$\star\star\star$}\label{deep:ex:ode_source_deep:blowup}\textbf{题目} 解$y'=y^2,y(0)=a$并给最大区间。\par\textbf{思路} 显式解的极点是有限时爆破。\par\textbf{详细解答} $a=0$为零解；$a\ne0$时$y=\frac{a}{1-at}$。若$a>0$，含0的最大区间$(-\infty,\frac{1}{a})$；若$a<0$则$(\frac{1}{a},\infty)$。在有限端点$|y|\to\infty$，不能连续延拓。若解始终位于$f$定义域的某个紧集内，局部存在定理可在端点重复应用，这说明本例不能延拓的原因。\par\textbf{方法点评} 局部唯一不排除有限时间爆破。\par\textbf{变式与易错点} 写出有理表达式不意味着可跨越极点连接为一个解。\end{example}

\begin{theorem}{Gronwall不等式与连续依赖}若$u$在$[0,T]$上非负连续，且对$0\le t\le T$有$u(t)\le A+B\int_0^tu(s)\dd s$，$A,B\ge0$，则$u(t)\le A\e^{Bt}$。\end{theorem}\begin{proof}令$v(t)=A+B\int_0^tu$，则$u\le v,v'\le Bv$，故$(\e^{-Bt}v)'\le0$，得到$v\le A\e^{Bt}$；$A=0$也可通过$A+\varepsilon$再极限处理。两解差满足此式时得到$|y(t)-z(t)|\le|y(0)-z(0)|\e^{Lt}$。\end{proof}

\begin{example}{竞赛拓展｜初值可微与变分方程}{难度：$\star\star\star$}\label{deep:ex:ode_source_deep:sensitivity}\textbf{题目} 设$y'=f(t,y)$且$f_y$连续，说明$z=\frac{\partial y}{\partial a}$（$y(0,a)=a$）满足什么方程，并在$y'=y^2$核验。\par\textbf{思路} 对积分方程求参数导数，得到线性变分式。\par\textbf{详细解答} 在公共紧存在区间，差商通过连续依赖与$f_y$的一致连续性收敛，满足$z'=f_y(t,y(t,a))z,z(0)=1$，所以$z=\exp(\int_0^tf_y(s,y(s,a))\dd s)$。对$y=\frac{a}{1-at}$，直接算$z=(1-at)^{-2}$；变分式$z'=2a(1-at)^{-1}z$给同值。\par\textbf{方法点评} 初值依赖可通过线性方程计算，即使原方程非线性。\par\textbf{变式与易错点} 必须在共同存在区间内讨论，不能越过参数相关爆破时间。\end{example}

\subsection{Wronskian、降阶与非齐次结构}\label{deep:ode_source_deep:linear}

\begin{method}{Wronskian、降阶与非齐次结构}\textbf{适用条件} 线性方程最高阶系数不为0，归一化后系数连续。\par\textbf{步骤与依据} 先求齐次基，再取一个非齐次特解；已知一解可降阶；常数变易用Wronskian解线性方程组。\par\textbf{易错点} 线性无关解数目应等于阶数；不同区间的奇点需分开。\end{method}

\begin{theorem}{Abel恒等式与二阶降阶}$y''+p(x)y'+q(x)y=0$的两解Wronskian$W=y_1y_2'-y_1'y_2$满足$W'=-pW$。若$y_1\ne0$，第二解可取$y_2=y_1\int\frac{\e^{-\int p}}{y_1^2}\dd x$。\end{theorem}\begin{proof}微分$W$后代入方程，$q$项抵消。另设$y_2=vy_1$，代入得$y_1v''+(2y_1'+py_1)v'=0$，一阶积分给$v'=\frac{\e^{-\int p}}{y_1^2}$，积分即公式。\end{proof}

\begin{example}{考研风格｜已知解降阶与Euler双解}{难度：$\star\star\star$}\label{deep:ex:ode_source_deep:reduction}\textbf{题目} 已知$y_1=x$，解$x^2y''-2xy'+2y=0$，$x>0$。\par\textbf{思路} 用降阶公式或Euler幂函数试解。\par\textbf{详细解答} 归一化$p=-\frac{2}{x}$，$\e^{-\int p}=x^2$，故$y_2=x\int \frac{x^2}{x^2}\dd x=x^2$，通解$C_1x+C_2x^2$。第二法设$y=x^r$，指标式$r(r-1)-2r+2=(r-1)(r-2)=0$，同解。\par\textbf{方法点评} 降阶公式适合已知一解；Euler法更快但依赖等幂结构。\par\textbf{变式与易错点} 若$x<0$可在该半轴独立处理，不能跨0套归一化。\end{example}

\begin{example}{考研风格｜常数变易与卷积双解}{难度：$\star\star\star$}\label{deep:ex:ode_source_deep:variation}\textbf{题目} 解$y''+y=\sec x$，在$\cos x\ne0$的区间。\par\textbf{思路} 齐次基$\cos x,\sin x$，Wronskian为1。\par\textbf{详细解答} 常数变易给$u_1'=-\sin x\sec x=-\tan x,u_2'=\cos x\sec x=1$，可取$u_1=\ln|\cos x|,u_2=x$。通解$y=C_1\cos x+C_2\sin x+\cos x\ln|\cos x|+x\sin x$。第二法以Green核写特解$y_p=\int_{x_0}^x\sin(x-t)\sec t\dd t$，展开$\sin(x-t)$并积分，差异仅为齐次项。\par\textbf{方法点评} 卷积核形式天然适合初值，常数变易便于求闭式。\par\textbf{变式与易错点} 非齐次项不能用有限多项式待定系数法。\end{example}

\subsection{微分算子移位、共振与复指数}\label{deep:ode_source_deep:operator}

\begin{method}{微分算子移位、共振与复指数}\textbf{适用条件} 常系数线性算子$P(D)$作用于指数多项式或三角函数。\par\textbf{步骤与依据} 用$P(D)[\e^{ax}v]=\e^{ax}P(D+a)v$；共振时提取$D^m$，再积分$m$次。\par\textbf{易错点} 形式逆算子必须说明作用函数族，不能把$\frac{1}{D}$当普通数。\end{method}

\begin{theorem}{移位公式}对多项式$P$，$P(D)(\e^{ax}v)=\e^{ax}P(D+a)v$。若$P(\lambda)$在$a$有$m$重根，则$P(D+a)=D^mQ(D)$且$Q(0)\ne0$。\end{theorem}\begin{proof}基本式$D(\e^{ax}v)=\e^{ax}(D+a)v$，归纳得$D^k$的对应式，线性相加即可。重根因式分解将共振转为积分因子$x^m$，说明待定系数为何须乘$x^m$。\end{proof}

\begin{example}{考研风格｜共振三解}{难度：$\star\star\star$}\label{deep:ex:ode_source_deep:resonance}\textbf{题目} 求$y''-2y'+y=\e^x$的通解。\par\textbf{思路} 特征根1为二重根。\par\textbf{详细解答} \textbf{待定系数法：}$y_p=Ax^2\e^x$，代入得$2A\e^x=\e^x$，所以$A=\frac{1}{2}$。\par\textbf{算子移位：}设$y=\e^xv$，$(D-1)^2y=\e^xv''$，故$v''=1$，$v=\frac{x^2}{2}+C_1x+C_2$。\par\textbf{连续一阶法：}设$z=y'-y$，先解$z'-z=\e^x$得$z=\e^x(x+C_1)$，再解$y'-y=z$得同解。\par\textbf{方法点评} 三法结果$y=\e^x(C_2+C_1x+\frac{x^2}{2})$；移位法在重根问题最简。\par\textbf{变式与易错点} 设$A\e^x$或$Ax\e^x$均落入齐次解，必须乘$x^2$。\end{example}

\begin{example}{考研风格｜复指数与实待定系数}{难度：$\star\star\star$}\label{deep:ex:ode_source_deep:complexforce}\textbf{题目} 求$y''+y=\e^x\cos x$的一个特解。\par\textbf{思路} 先解复方程右端$\e^{(1+i)x}$，再取实部。\par\textbf{详细解答} $P(1+i)=(1+i)^2+1=1+2i$，故复特解$\frac{\e^{(1+i)x}}{1+2i}$。取实部得$y_p=\frac{\e^x(\cos x+2\sin x)}{5}$。第二法设$\e^x(A\cos x+B\sin x)$，代入比较得$A=\frac{1}{5},B=\frac{2}{5}$。\par\textbf{方法点评} 复指数将两个实系数方程压缩为一个复数除法。\par\textbf{变式与易错点} 必须检查$P(1+i)\ne0$，否则要按重根次数乘$x$。\end{example}

\subsection{幂级数、生成函数与积分方程}\label{deep:ode_source_deep:series}

\begin{method}{幂级数、生成函数与积分方程}\textbf{适用条件} 系数解析，或题设给出可化为导数的递推和积分关系。\par\textbf{步骤与依据} 统一幂次后比较系数；积分方程求导时同时保留由端点产生的初值。\par\textbf{易错点} 求导可能丢积分常数，初值必须由原式恢复。\end{method}

\begin{example}{竞赛风格｜幂级数解与显式解互证}{难度：$\star\star\star$}\label{deep:ex:ode_source_deep:power}\textbf{题目} 用幂级数解$y''=xy$，$y(0)=1,y'(0)=0$。\par\textbf{思路} 把$x^n$系数对齐，按模3分组。\par\textbf{详细解答} 设$y=\sum a_nx^n$，则$a_2=0$，$(n+2)(n+1)a_{n+2}=a_{n-1}$（$n\ge1$）。由$a_0=1,a_1=0$得只有$a_{3k}$非零，$a_{3k}=\frac{1}{\prod_{j=1}^k(3j)(3j-1)}$。所以$y=1+\frac{x^3}{6}+\frac{x^6}{180}+\frac{x^9}{12960}+\cdots$，比值给任意有限$x$绝对收敛。另一独立解从$a_1=1,a_0=0$生成，它们构成Airy方程的基。\par\textbf{方法点评} 不需要把答案强行写成命名特殊函数；递推与初值已经唯一确定解。\par\textbf{变式与易错点} 比较系数时$x$因子使下标移动一位，$n=0$需单列。\end{example}

\begin{example}{竞赛风格｜生成函数型递推}{难度：$\star\star\star$}\label{deep:ex:ode_source_deep:generating}\textbf{题目} 设$a_0=4,a_1=1$，且$n(n-1)a_n=a_{n-2}$（$n\ge2$），求$S(x)=\sum_{n\ge0}a_nx^n$。\par\textbf{思路} 递推分母正好对应二阶导数。\par\textbf{详细解答} 逐项关系给$S''=S$，$S(0)=4,S'(0)=1$，所以$S=4\cosh x+\sinh x=\frac{5\e^x+3\e^{-x}}{2}$。直接分奇偶也可得$a_{2k}=\frac{4}{(2k)!}$、$a_{2k+1}=\frac{1}{(2k+1)!}$，故全实轴收敛并得到相同和。\par\textbf{方法点评} 微分方程法与奇偶系数法互证，避免先假定收敛却未解释。\par\textbf{变式与易错点} 递推关系必须对所有规定指标成立；初始系数决定微分方程的初值。\end{example}

\begin{example}{考研风格｜Volterra方程转ODE双解}{难度：$\star\star\star$}\label{deep:ex:ode_source_deep:volterra}\textbf{题目} 解$y(x)=1+\int_0^x(x-t)y(t)\dd t$。\par\textbf{思路} 连续求导两次，并从原式读取初值。\par\textbf{详细解答} 若$y$连续，右边先得$y'(x)=\int_0^xy(t)\dd t$，再得$y''=y$，初值$y(0)=1,y'(0)=0$，故$y=\cosh x$。代回积分$\int_0^x(x-t)\cosh t\dd t=\cosh x-1$验证。第二法迭代核得到$1+\frac{x^2}{2!}+\frac{x^4}{4!}+\cdots=\cosh x$。\par\textbf{方法点评} 微分法短，逐次逼近法直接显示积分算子结构。\par\textbf{变式与易错点} 只写$y''=y$而漏初值会多出两参数。\end{example}

\begin{example}{竞赛拓展｜Euler折线的误差分析}{难度：$\star\star\star$}\label{deep:ex:ode_source_deep:eulerscheme}\textbf{题目} 用显式Euler法近似$y'=y,y(0)=1$在$t=1$的值，步长$h=\frac{1}{n}$，并说明误差阶。\par\textbf{思路} 递推可显式求解，再用对数展开。\par\textbf{详细解答} $y_{k+1}=(1+h)y_k$，故$y_n=(1+\frac{1}{n})^n$。由对数展开$y_n=\e[1-\frac{1}{2n}+O(n^{-2})]$，所以全局误差$\e-y_n=\frac{\e h}{2}+O(h^2)$。每步精确流$\e^h=1+h+\frac{h^2}{2}+O(h^3)$，单步截断为$O(h^2)$，累积$n$步成为$O(h)$；一般情形再用离散Gronwall控制传播。\par\textbf{方法点评} 线素场与Euler折线不仅给出图形近似，还可用于量化误差。\par\textbf{变式与易错点} 局部截断二阶不意味着全局二阶。\end{example}

\subsection{周期系数与解的结构}\label{deep:ode_source_deep:periodic}

\begin{method}{周期系数与解的结构}\textbf{适用条件} 一阶线性系数连续且具有周期$T>0$。\par\textbf{步骤与依据} 将一个周期内的积分看作乘子；非齐次解之差归于齐次问题。\par\textbf{易错点} 系数周期不等于所有解自动具有同一周期。\end{method}

\begin{example}{竞赛风格｜周期解的充要条件}{难度：$\star\star\star$}\label{deep:ex:ode_source_deep:periodicA}\textbf{题目} $p$是实连续$T$周期函数。证明$y'+p(x)y=0$的非零解具有周期$T$，当且仅当$\int_0^Tp(t)\dd t=0$。\par\textbf{思路} 用显式解比较$y(x+T)$与$y(x)$。\par\textbf{详细解答} 非零解$y(x)=C\exp(-\int_0^xp(t)\dd t)$，$C\ne0$。由于$p$周期，$\int_x^{x+T}p(t)\dd t$对$x$的导数为$p(x+T)-p(x)=0$，故等于$\int_0^Tp$。于是$\frac{y(x+T)}{y(x)}=\exp(-\int_0^Tp)$。实指数等于1当且仅当指数为0，充要条件得证。另一法直接积分$\frac{y'}{y}=-p$，周期性迫使$\ln|y(x+T)|-\ln|y(x)|=0$，得到必要性；充分性由同一积分还原。\par\textbf{方法点评} 周期积分与单周期乘子把“系数周期”转为“解周期”的准确条件。\par\textbf{变式与易错点} 零解无论积分是否为0都周期，所以题目必须限定非零解；复系数时指数为1还允许$2\pi i$的整数倍。\end{example}

\begin{example}{考研风格｜两个非齐次特解生成通解}{难度：$\star\star\star$}\label{deep:ex:ode_source_deep:twosolutions}\textbf{题目} 若$y_1,y_2$是$y'+p(x)y=q(x)$的两个不同解，写出通解并说明任意线性组合何时仍是解。\par\textbf{思路} 解的差满足齐次方程。\par\textbf{详细解答} $h=y_1-y_2$满足$h'+ph=0$，且不是恒零，所以一阶齐次通解为$Ch$。原方程通解为$y=y_2+C(y_1-y_2)$。一般$\alpha y_1+\beta y_2$代入左侧得到$(\alpha+\beta)q$，故当$q$不恒零时，$\alpha+\beta=1$才保证同为非齐次解；$q\equiv0$时任意线性组合均可。\par\textbf{方法点评} 线性结构说明通解为何可写成“齐通加非特”。\par\textbf{变式与易错点} “两个不同解”在唯一性条件下意味着差在整个区间不为0，不能选相同解作基。\end{example}

\section{方程组、稳定性、差分与一阶偏微分}\label{deep:ode_qualitative_deep}

\begin{examtip}本节讨论方程组的结构和定性结论，并拓展差分与特征线内容。\end{examtip}

\subsection{基本矩阵、Jordan块与常数变易}\label{deep:ode_qualitative_deep:systems}

\begin{method}{基本矩阵、Jordan块与常数变易}\textbf{适用条件} $X'=A(t)X+b(t)$，$A,b$连续。\par\textbf{步骤与依据} 基本矩阵$\Phi$可逆；设$X=\Phi C$，得到$C'=\Phi^{-1}b$。常系数时计算$\e^{At}$。\par\textbf{易错点} 重复特征值未必有足够特征向量，要检查Jordan链。\end{method}

\begin{theorem}{矩阵Liouville公式}若$\Phi'=A\Phi,\Phi(t_0)$可逆，则$\det\Phi(t)=\det\Phi(t_0)\exp(\int_{t_0}^t\operatorname{tr}A(s)\dd s)$。\end{theorem}\begin{proof}在行列式非零的局部区间，由Jacobi公式$(\det\Phi)'=\det\Phi\,\operatorname{tr}(\Phi^{-1}\Phi')=\operatorname{tr}A\det\Phi$。解一阶方程得到非零表达，故不能在有限端点首次变零，从而延伸到整个存在区间。\end{proof}

\begin{example}{考研风格｜Jordan与逐次消元双解}{难度：$\star\star\star$}\label{deep:ex:ode_qualitative_deep:jordan}\textbf{题目} 解$x'=x+y,y'=y$。\par\textbf{思路} 系数矩阵为$I+N,N^2=0$。\par\textbf{详细解答} 矩阵指数$\e^{At}=\e^t(I+tN)$，所以$(x,y)^T=\e^t(C_1+tC_2,C_2)^T$。第二法先得$y=C_2\e^t$，再解$x'-x=C_2\e^t$，积分得$x=\e^t(C_1+C_2t)$。\par\textbf{方法点评} 多项式$t$是缺少特征向量时Jordan结构的信号。\par\textbf{变式与易错点} 只有一个特征值不代表解只有一个自由常数。\end{example}

\begin{example}{竞赛风格｜复特征值与极坐标}{难度：$\star\star\star$}\label{deep:ex:ode_qualitative_deep:rotation}\textbf{题目} 解$x'=-x-y,y'=x-y$并判断原点稳定性。\par\textbf{思路} 复变量$z=x+iy$，或直接求$r,\theta$。\par\textbf{详细解答} $z'=(-1+i)z$，故$z=z_0\e^{-t}\e^{it}$，轨道为旋转向内的螺线。极坐标法$r'=\frac{xx'+yy'}{r}=-r$，$\theta'=\frac{xy'-yx'}{r^2}=1$，同得$r=r_0\e^{-t},\theta=t+\theta_0$。原点全局渐近稳定。\par\textbf{方法点评} 复数法给显式解，极坐标法解释相图。\par\textbf{变式与易错点} $\theta$在原点无定义，原点平衡解需单列。\end{example}

\subsection{Lyapunov函数、线性化与极限环}\label{deep:ode_qualitative_deep:stability}

\begin{method}{Lyapunov函数、线性化与极限环}\textbf{适用条件} 自治系统的平衡点附近，向量场至少连续可微。\par\textbf{步骤与依据} 构造正定$V$并计算沿轨道导数；严格负定给渐近稳定。非双曲点须看高阶项。\par\textbf{易错点} 线性化有零实部特征值时结论通常不定。\end{method}

\begin{example}{竞赛拓展｜线性化失效的两例}{难度：$\star\star\star$}\label{deep:ex:ode_qualitative_deep:nonhyperbolic}\textbf{题目} 比较$x'=-x^3$与$x'=x^3$在0的稳定性。\par\textbf{思路} 两者线性化都为$x'=0$，用$V=\frac{x^2}{2}$区分。\par\textbf{详细解答} 第一式$\dot V=-x^4<0$（$x\ne0$），显式解$x(t)=\frac{x_0}{\sqrt{1+2x_0^2t}}$，趋0，故渐近稳定。第二式$\dot V=x^4>0$，解$x=\frac{x_0}{\sqrt{1-2x_0^2t}}$，非零初值远离0且有限时爆破，所以不稳定。\par\textbf{方法点评} 同一零特征值可以对应相反的非线性行为。\par\textbf{变式与易错点} 不能把“线性化不渐近稳定”当作原系统不稳定。\end{example}

\begin{example}{竞赛拓展｜可计算的稳定极限环}{难度：$\star\star\star$}\label{deep:ex:ode_qualitative_deep:limitcycle}\textbf{题目} 研究$x'=x(1-x^2-y^2)-y,\ y'=y(1-x^2-y^2)+x$。\par\textbf{思路} 极坐标把径向与角向分离。\par\textbf{详细解答} $r'=r(1-r^2),\theta'=1$。$r=0$是平衡，$r=1$给周期$2\pi$的闭轨道。若$0<r_0<1$，$r$递增；若$r_0>1$，$r$递减，均趋1。闭轨道附近不存在其他闭轨道，因为任何$r\ne1$沿轨道严格单调，所以$r=1$是吸引极限环。原点因$r'>0$不稳定。\par\textbf{方法点评} 极限环需要孤立性；一般中心的一族圆周不是极限环。\par\textbf{变式与易错点} 不能仅看$\theta'=1$就断言每条轨道周期。\end{example}

\begin{theorem}{Bendixson排除闭轨道的一个版本}在平面单连通区域$D$，$C^1$场$F=(P,Q)$的散度$\partial_xP+\partial_yQ$处处严格同号，则$D$内没有简单闭轨道。\end{theorem}\begin{proof}若存在闭轨道$\Gamma$围成$\Omega\subset D$，沿轨道场为切向，外法向通量为0。由Green通量公式$0=\int_\Gamma F\cdot n\dd s=\iint_\Omega\operatorname{div}F\dd A$，与严格同号矛盾。单连通性保证内部不穿出$D$。\end{proof}

\subsection{线性差分、共振与离散稳定性}\label{deep:ode_qualitative_deep:difference}

\begin{method}{线性差分、共振与离散稳定性}\textbf{适用条件} 递推的系数为常数或能用变量替换线性化。\par\textbf{步骤与依据} 齐次根判增长；外力与特征根重合时乘$n$；稳定性比较特征根模是否小于1。\par\textbf{易错点} 连续时间看实部，离散时间看模长。\end{method}

\begin{example}{竞赛拓展｜离散共振与生成函数}{难度：$\star\star\star$}\label{deep:ex:ode_qualitative_deep:diffresonance}\textbf{题目} 解$a_{n+1}-2a_n=2^n,a_0=0$。\par\textbf{思路} 除以$2^{n+1}$化为等差增量。\par\textbf{详细解答} 令$b_n=\frac{a_n}{2^n}$，则$b_{n+1}-b_n=\frac{1}{2}$，所以$a_n=n2^{n-1}$。生成函数法$A(z)=\sum a_nz^n$满足$\frac{A}{z}-2A=\frac{1}{1-2z}$，故$A=\frac{z}{(1-2z)^2}$，展开系数也为$n2^{n-1}$。\par\textbf{方法点评} $n$因子是离散共振，与ODE乘$x$完全对应。\par\textbf{变式与易错点} 生成函数等式先在$|z|<\frac{1}{2}$理解，也可在形式幂级数中操作。\end{example}

\begin{example}{考研风格｜二阶差分与稳定性}{难度：$\star\star\star$}\label{deep:ex:ode_qualitative_deep:secondrec}\textbf{题目} 解$a_{n+2}-\frac{a_{n+1}}{2}-\frac{a_n}{2}=0,a_0=0,a_1=1$。\par\textbf{思路} 特征根为1与$-\frac{1}{2}$。\par\textbf{详细解答} $a_n=C+D(-\frac{1}{2})^n$，初值得$C=\frac{2}{3},D=-\frac{2}{3}$，故$a_n=\frac23[1-(-\frac{1}{2})^n]\to\frac{2}{3}$。若改变初值，单位根1保留一维不衰减分量，因此零解稳定但不渐近稳定；$|-\frac{1}{2}|<1$的分量交替衰减。\par\textbf{方法点评} 有界性与趋零不同，由单位圆上的根决定。\par\textbf{变式与易错点} 若单位根带非平凡Jordan块，会出现$n$增长，稳定性又不同。\end{example}

\subsection{一阶偏微分方程的特征线}\label{deep:ode_qualitative_deep:pde}

\begin{method}{一阶偏微分方程的特征线}\textbf{适用条件} 一阶线性或拟线性方程，初始曲线不与特征方向相切。\par\textbf{步骤与依据} 写$\dot x=a,\dot y=b,\dot u=c$；找第一积分，再用初始条件确定任意函数。\par\textbf{易错点} 初始曲线特征时可能无解或不唯一，不能照常解参数。\end{method}

\begin{example}{竞赛拓展｜输运方程的两解}{难度：$\star\star\star$}\label{deep:ex:ode_qualitative_deep:transport}\textbf{题目} 解$u_t+cu_x=0,u(x,0)=g(x)$。\par\textbf{思路} 沿$x-ct$不变的直线传播。\par\textbf{详细解答} 特征式$\dot t=1,\dot x=c,\dot u=0$，所以$u(x,t)=g(x-ct)$。第二法换坐标$\xi=x-ct,\tau=t$，链式法则给$u_t+cu_x=v_\tau$，故$v$与$\tau$无关，由初值得$v=g(\xi)$。\par\textbf{方法点评} 特征线法描述传播，坐标法直接解释算子的简化。\par\textbf{变式与易错点} 若$g$不光滑，公式仍可给弱解，但不能无条件称为经典解。\end{example}

\begin{example}{竞赛拓展｜非线性特征与破裂时间}{难度：$\star\star\star$}\label{deep:ex:ode_qualitative_deep:quasilinear}\textbf{题目} 求$u_t+uu_x=0,u(x,0)=g(x)$的局部隐式解。\par\textbf{思路} 沿特征$u$为常数，但速度由初值决定。\par\textbf{详细解答} 以初始点$\xi$参数化，$u=g(\xi),x=\xi+t g(\xi)$。当$1+tg'(\xi)\ne0$时可局部反解$\xi$，得$u(x,t)=g(x-tu(x,t))$。进一步$u_x=\frac{g'(\xi)}{1+tg'(\xi)}$。若$\inf g'<0$，正向最早梯度爆破候选为$T=-\frac{1}{\inf g'}$（下确界达到且适当正则时），特征相交后需要弱解与熵条件。\par\textbf{方法点评} 同样的特征方法展示经典解为何只能局部存在。\par\textbf{变式与易错点} 在特征相交后不能保留多值函数作经典解。\end{example}

\section{微分方程的换元与退化分支}\label{ic:delta_ode}

\begin{examtip}本节逐步演示倒置自变量、仿射齐次化、高阶Euler方程及非常系数坐标变换。\end{examtip}

\subsection{将自变量与因变量交换}\label{ic:delta_ode:invert}\begin{method}{将自变量与因变量交换}\textbf{识别信号} $y'(x)$的倒数作为$x(y)$方程更简单。\par \textbf{操作步骤} 先单列$y'=0$的候选支路，再在可逆段写$\frac{\dd x}{\dd y}=\frac{1}{y'}$，解新方程，最后代回原式。\par \textbf{疑难点与条件} 不是全局每个解都可反函数化；常值解可能被倒置丢掉。\end{method}

\begin{example}{竞赛风格｜倒置后变一阶线性}{难度：$\star\star\star$}\label{ic:ex:delta_ode:invertcase}\textbf{题目} 解$y'=\frac{y}{2x+y}$。\par \textbf{思路} 非零支路看作$x(y)$。\par \textbf{详细解答} 当$y\ne0$时，$\frac{\dd x}{\dd y}-\frac{2x}{y}=1$，积分因子$y^{-2}$给$(\frac{x}{y^2})'=y^{-2}$，所以$x=Cy^2-y$，即$x+y=Cy^2$。另有$y\equiv0$，它在不含$x=0$的区间满足原式。\par \textbf{验算与条件复核} 隐式式求导得$1+y'=2Cyy'$，结合$2Cy=\frac{2x+2y}{y}$可还原原方程；还须排除$2x+y=0$。\par \textbf{疑难点与方法比较} 补零解并检查原分母，是倒置法的两个关键。\end{example}

\subsection{仿射分式的平移齐次化}\label{ic:delta_ode:affine}\begin{method}{仿射分式的平移齐次化}\textbf{识别信号} $y'=\frac{a_1x+b_1y+c_1}{a_2x+b_2y+c_2}$。\par \textbf{操作步骤} 若两行线性系数不成比例，先解两条分子分母零线的交点，再平移消常数并设$v=\frac{Y}{X}$；若成比例，取公共线性组合为变量，转成可分离式。\par \textbf{疑难点与条件} 换元中除掉$v$或$1-v$时单列特殊直线解；最终解可在原方程正则处延拓。\end{method}

\begin{example}{竞赛风格｜平移、比值与特殊解}{难度：$\star\star\star$}\label{ic:ex:delta_ode:affinecase}\textbf{题目} 解$y'=\frac{2y+4}{x+y-1}$。\par \textbf{思路} 两零线交于$(3,-2)$。\par \textbf{详细解答} 取$X=x-3,Y=y+2$，得$Y'=\frac{2Y}{X+Y}$。在$X\ne0$令$v=\frac{Y}{X}$，则$Xv'=\frac{v(1-v)}{1+v}$，分离积分$\ln|v|-2\ln|1-v|=\ln|X|+C$。还原得$(x-y-5)^2=C_1(y+2)$。特殊支$v=1$为$y=x-5$，已在$C_1=0$中；被除掉的$v=0$给额外解$y=-2$。\par \textbf{验算与条件复核} 对隐式式求导并用该式消$C_1$，得到$y'=\frac{2(y+2)}{x+y-1}$。所有解均只在原分母非零的正则区间取用。\par \textbf{疑难点与方法比较} 隐式常数改写可能丢常值支路，不能只给一个形式就停止检查。\end{example}

\begin{example}{竞赛风格｜三阶Euler的算子展开}{难度：$\star\star\star$}\label{ic:ex:delta_ode:eulerthirdorder}\textbf{题目} 解$x^3y'''+x^2y''-4xy'=3x^2$，$x>0$。\par \textbf{思路} 令$t=\ln x$，逐阶保留$D(D-1)$结构。\par \textbf{详细解答} 有$xy'=Dv,x^2y''=D(D-1)v,x^3y'''=D(D-1)(D-2)v$。方程化为$D(D-3)(D+1)v=3\e^{2t}$。齐次解$C_1+C_2\e^{-t}+C_3\e^{3t}$；因$P(2)=-6$，特解$-\frac{\e^{2t}}{2}$。所以$y=C_1+\frac{C_2}{x}+C_3x^3-\frac{x^2}{2}$。\par \textbf{验算与条件复核} 原算子作用于$x^r$给$r(r-3)(r+1)x^r$，对$r=2$给$-6x^2$，验证全部分量。\par \textbf{疑难点与方法比较} 不能把$x^3y'''$误写为单纯$D^3v$。\end{example}

\begin{example}{竞赛风格｜非常系数方程的坐标拉直}{难度：$\star\star\star$}\label{ic:ex:delta_ode:tangentcoordinate}\textbf{题目} 解$\cos^2x\,y''-(\sin2x-2)y'+\sec^2x\,y=\tan^3x+\tan x$，在$\cos x\ne0$区间。\par \textbf{思路} 设$t=\tan x,y=v(t)$，最高阶系数与$t'$相配。\par \textbf{详细解答} $y'=(1+t^2)v'$、$y''=(1+t^2)^2v''+2t(1+t^2)v'$。代入后全部方程除以$1+t^2$，得到$v''+2v'+v=t$。齐次解$(C_1+C_2t)\e^{-t}$，一次特解$t-2$。答案$y=(C_1+C_2\tan x)\e^{-\tan x}+\tan x-2$。\par \textbf{验算与条件复核} 链式导数代入显示$v'$中的$2t$项先相消，余下常系数式；$v=t-2$代入为$t$。\par \textbf{疑难点与方法比较} 必须计算二阶链式的附加项，不能只平方一阶变换系数。\end{example}

\section{图解：方向场、平衡解与增长}
\begin{figure}[htbp]\centering
\begin{tikzpicture}\begin{axis}[width=.8\textwidth,height=6cm,axis lines=middle,xmin=-2,xmax=4,ymin=-.15,ymax=1.5,xlabel=$x$,ylabel=$y$,samples=120]
\addplot[Teal,thick,domain=-2:4]{1/(1+3*exp(-x))};
\addplot[dashed,Ink,domain=-2:4]{1};\addplot[dashed,Ink,domain=-2:4]{0};
\draw[Gold] (axis cs:-1.6300,0.1792)--(axis cs:-1.3700,0.2208);
\draw[Gold] (axis cs:-0.6300,0.1792)--(axis cs:-0.3700,0.2208);
\draw[Gold] (axis cs:0.3700,0.1792)--(axis cs:0.6300,0.2208);
\draw[Gold] (axis cs:1.3700,0.1792)--(axis cs:1.6300,0.2208);
\draw[Gold] (axis cs:2.3700,0.1792)--(axis cs:2.6300,0.2208);
\draw[Gold] (axis cs:3.3700,0.1792)--(axis cs:3.6300,0.2208);
\draw[Gold] (axis cs:-1.6300,0.4675)--(axis cs:-1.3700,0.5325);
\draw[Gold] (axis cs:-0.6300,0.4675)--(axis cs:-0.3700,0.5325);
\draw[Gold] (axis cs:0.3700,0.4675)--(axis cs:0.6300,0.5325);
\draw[Gold] (axis cs:1.3700,0.4675)--(axis cs:1.6300,0.5325);
\draw[Gold] (axis cs:2.3700,0.4675)--(axis cs:2.6300,0.5325);
\draw[Gold] (axis cs:3.3700,0.4675)--(axis cs:3.6300,0.5325);
\draw[Gold] (axis cs:-1.6300,0.7792)--(axis cs:-1.3700,0.8208);
\draw[Gold] (axis cs:-0.6300,0.7792)--(axis cs:-0.3700,0.8208);
\draw[Gold] (axis cs:0.3700,0.7792)--(axis cs:0.6300,0.8208);
\draw[Gold] (axis cs:1.3700,0.7792)--(axis cs:1.6300,0.8208);
\draw[Gold] (axis cs:2.3700,0.7792)--(axis cs:2.6300,0.8208);
\draw[Gold] (axis cs:3.3700,0.7792)--(axis cs:3.6300,0.8208);
\draw[Gold] (axis cs:-1.6300,1.2312)--(axis cs:-1.3700,1.1688);
\draw[Gold] (axis cs:-0.6300,1.2312)--(axis cs:-0.3700,1.1688);
\draw[Gold] (axis cs:0.3700,1.2312)--(axis cs:0.6300,1.1688);
\draw[Gold] (axis cs:1.3700,1.2312)--(axis cs:1.6300,1.1688);
\draw[Gold] (axis cs:2.3700,1.2312)--(axis cs:2.6300,1.1688);
\draw[Gold] (axis cs:3.3700,1.2312)--(axis cs:3.6300,1.1688);
\draw[Gold] (axis cs:-1.6300,1.4728)--(axis cs:-1.3700,1.3272);
\draw[Gold] (axis cs:-0.6300,1.4728)--(axis cs:-0.3700,1.3272);
\draw[Gold] (axis cs:0.3700,1.4728)--(axis cs:0.6300,1.3272);
\draw[Gold] (axis cs:1.3700,1.4728)--(axis cs:1.6300,1.3272);
\draw[Gold] (axis cs:2.3700,1.4728)--(axis cs:2.6300,1.3272);
\draw[Gold] (axis cs:3.3700,1.4728)--(axis cs:3.6300,1.3272);
\node[anchor=west]at(axis cs:2.5,1.12){$y=1$};
\node[anchor=west]at(axis cs:.2,.24){$y(0)=\frac{1}{4}$};
\end{axis}\end{tikzpicture}
\caption{$y'=y(1-y)$的方向场。$0<y<1$时斜率正，$y>1$时斜率负；$y=0,1$是平衡解。实线是$y=\frac{1}{1+3\e^{-x}}$，它始终在$(0,1)$内并趋向1。短线段只表示斜率方向，未表示固定时间长度。}\label{fig:review:ode}
\end{figure}

\section{数三应用：一阶常系数线性差分方程}\label{r2:difference}
离散变量以 $t=0,1,2,\ldots$ 编号。定义 $\Delta x_t=x_{t+1}-x_t$，不要把 $\Delta x_t$ 当作连续导数。方程 $\Delta x_t+\alpha x_t=b$ 等价于 $x_{t+1}=a x_t+b$，其中 $a=1-\alpha$。

\begin{theorem}{通解、特解与初值}
对常数 $a,b$ 及给定 $x_0$，递推 $x_{t+1}=a x_t+b$ 的唯一解为
\[
 x_t=\begin{cases}
 \displaystyle\frac{b}{1-a}+a^t\left(x_0-\frac{b}{1-a}\right),&a\ne1,\\[4pt]
 x_0+bt,&a=1.
 \end{cases}
\]
当 $a=0$ 时，直接规定 $t=0$ 用初值、$t\ge1$ 为 $b$，无需讨论符号 $0^0$。$a\ne1$ 时常数特解 $x_*=b/(1-a)$，齐次解是 $Ca^t$；$a=1$ 时常数特解失效，改试 $bt$。
\end{theorem}
\begin{proof}
当 $a\ne1$，令 $u_t=x_t-x_*$，则 $u_{t+1}=au_t$，逐次递推得到 $u_t=a^tu_0$；$a=1$ 时逐期差为 $b$，累加即可。每期由上一期唯一确定，故满足初值的解唯一。
\end{proof}
更一般的 $x_{t+1}=a x_t+b_t$ 可写成
\[
 x_t=a^t x_0+\sum_{j=0}^{t-1}a^{t-1-j}b_j\quad(t\ge1).
\]
对指数项 $b_t=c r^t$，若 $r\ne a$ 可试 $K r^t$ 并得 $K=c/(r-a)$；共振 $r=a\ne0$ 时改试 $Kt a^t$。零值退化情形直接用递推或有限和处理。

\subsection{平衡与收敛：检查系数模长}
当 $|a|<1$，所有初值收敛到 $x_*$。$0<a<1$ 时偏离同号衰减，$-1<a<0$ 时交替衰减。$a=-1$ 时非平衡初值一般为二周期；$|a|>1$ 时非平衡偏离的模增长，但初值恰为平衡仍是常数解。$a=1,b=0$ 时所有初值都平衡，$b\ne0$ 时没有常数平衡。

\begin{example}{常量外生支出与收敛路径}{数三应用}
$Y_{t+1}=0.8Y_t+10$，$Y_0=0$。平衡为 $50$，所以 $Y_t=50(1-0.8^t)$。核对 $Y_1=10,Y_2=18$；误差为 $50\cdot0.8^t$，与“每期增加0.8”无关。
\end{example}

\subsection{蛛网模型：滞后供给与市场出清}
设需求 $D_t=\alpha-\beta p_t$、供给 $S_t=\gamma+\delta p_{t-1}$，$\beta,\delta>0$。每期 $D_t=S_t$，得到
\[
 p_t=\frac{\alpha-\gamma}{\beta}-\frac{\delta}{\beta}p_{t-1},\qquad
 p_* =\frac{\alpha-\gamma}{\beta+\delta},\qquad
 p_t-p_* =\left(-\frac{\delta}{\beta}\right)^t(p_0-p_*).
\]
因此局部偏离收敛要求 $\delta<\beta$。这是按需求和供给都写成价格函数的斜率约定得到的条件，不应脱离坐标方向背诵“哪条曲线更陡”。价格、数量非负的经济可行性须另查。
例如 $\alpha=100,\gamma=10,\beta=2,\delta=1,p_0=20$，则 $p_*=30,p_1=35,p_2=27.5$，逐期振荡衰减。

\paragraph{自测与解答}
\begin{enumerate}
\item $x_{t+1}=-x_t/2+3,x_0=0$。答：$x_t=2-2(-1/2)^t$，前两期为 $3,3/2$，极限为 $2$。
\item $x_{t+1}=2x_t+3\cdot2^t,x_0=1$。答：共振特解为 $(3/2)t2^t$，所以 $x_t=2^t(1+3t/2)$；代 $t=0,1$ 核对为 $1,5$。
\item $x_{t+1}=x_t+2,x_0=1$。答：$x_t=1+2t$，本题不能使用 $b/(1-a)$。
\end{enumerate}


\section*{章末分层习题}\addcontentsline{toc}{section}{章末分层习题}
\refstepcounter{exercise}\label{q:4:1}\par\noindent\textbf{习题\theexercise\quad 基础题} 解$y\prime=3y,y(0)=2$。\quad\hyperref[a:4:1]{[答案]}\par
\refstepcounter{exercise}\label{q:4:2}\par\noindent\textbf{习题\theexercise\quad 基础题} 解$y\prime+x=0$。\quad\hyperref[a:4:2]{[答案]}\par
\refstepcounter{exercise}\label{q:4:3}\par\noindent\textbf{习题\theexercise\quad 基础题} 解$y\prime+y=1$。\quad\hyperref[a:4:3]{[答案]}\par
\refstepcounter{exercise}\label{q:4:4}\par\noindent\textbf{习题\theexercise\quad 考研风格} 解$y\prime\prime+4y=0$。\quad\hyperref[a:4:4]{[答案]}\par
\refstepcounter{exercise}\label{q:4:5}\par\noindent\textbf{习题\theexercise\quad 考研风格} 解$y\prime\prime-2y\prime+y=0$。\quad\hyperref[a:4:5]{[答案]}\par
\refstepcounter{exercise}\label{q:4:6}\par\noindent\textbf{习题\theexercise\quad 考研风格} 求$y\prime\prime+y=x$的通解。\quad\hyperref[a:4:6]{[答案]}\par
\refstepcounter{exercise}\label{q:4:7}\par\noindent\textbf{习题\theexercise\quad 竞赛拓展} 解$x^2y\prime\prime+xy\prime-y=0$，$x>0$。\quad\hyperref[a:4:7]{[答案]}\par
\refstepcounter{exercise}\label{q:4:8}\par\noindent\textbf{习题\theexercise\quad 竞赛拓展} 解$y(x)=2+\int_0^xy(t)\dd t$。\quad\hyperref[a:4:8]{[答案]}\par
\refstepcounter{exercise}\label{q:4:9}\par\noindent\textbf{习题\theexercise\quad 竞赛拓展} 求$y\prime\prime+y=0$且$y(0)=0,y\prime(0)=1$的幂级数。\quad\hyperref[a:4:9]{[答案]}\par
\section*{章末方法速查}\addcontentsline{toc}{section}{章末方法速查}
\par\noindent\begin{tabularx}{\textwidth}{p{23mm}XXX}\toprule 方法&适用条件&典型信号&易错点\\\midrule
分离变量&可分积&$p(x)q(y)$&漏平衡解\\积分因子&一阶线性&$y\prime+py$&未化标准型\\待定系数&导数封闭族&指数多项式&漏共振因子\\常数变易&基解独立&一般激励&Wronskian符号\\
\bottomrule\end{tabularx}
